% Generated draft: requires mathematical and source review.
\documentclass[11pt]{article}
\usepackage[margin=25mm]{geometry}
\usepackage{fontspec}
\setmainfont{Latin Modern Roman}
\usepackage{amsmath,amssymb,mathtools,unicode-math}
\setmathfont{Latin Modern Math}
\usepackage{enumitem,needspace,xurl,hyperref}
\setlength{\parindent}{0pt}
\setlength{\parskip}{5pt}
\setlength{\emergencystretch}{3em}
\newcommand{\N}{\mathbb{N}}
\newcommand{\Z}{\mathbb{Z}}
\newcommand{\Q}{\mathbb{Q}}
\newcommand{\R}{\mathbb{R}}
\newcommand{\C}{\mathbb{C}}
\newcommand{\F}{\mathbb{F}}
\newcommand{\A}{\mathbb{A}}
\newcommand{\PP}{\mathbb{P}}
\newcommand{\E}{\mathbb{E}}
\newcommand{\eps}{\varepsilon}
\newcommand{\dd}{\,\mathrm{d}}
\begin{document}
\section*{UnsolvedMath ID: 30001576}
\textbf{Source problem code:} catalogue entry
\subsection{Definitions and mathematical statement}
Under the conventional complex-analytic interpretation, let \(g\geq 1\) be an integer and let \(\mathcal A_g\) denote the moduli space of complex principally polarized abelian varieties of dimension \(g\). The source record does not explicitly specify the base field or the range of genera; these conventions therefore require confirmation.

For an analytic description, put
\[
\mathfrak H_g=\{\tau\in M_g(\C):\tau^{\mathsf T}=\tau,\ \operatorname{Im}\tau>0\},
\qquad
A_\tau=\C^g/(\Z^g+\tau\Z^g).
\]
For \(\varepsilon,\delta\in\{0,1\}^g\), define the theta function with characteristic by
\[
\theta\!\begin{bmatrix}\varepsilon\\\delta\end{bmatrix}(\tau,z)
=
\sum_{n\in\Z^g}
\exp\!\left(
\pi i(n+\varepsilon/2)^{\mathsf T}\tau(n+\varepsilon/2)
+2\pi i(n+\varepsilon/2)^{\mathsf T}(z+\delta/2)
\right).
\]
The characteristic is called odd if \(\varepsilon^{\mathsf T}\delta\equiv1\pmod 2\). Relative to the standard symmetric theta divisor on \(A_\tau\), it specifies the odd two-torsion point represented by \((\tau\varepsilon+\delta)/2\). Here a two-torsion point is a point \(x\in A_\tau\) satisfying \(2x=0\). The multiplicity of a theta divisor at a point means the vanishing order of a local defining equation there; in these coordinates it equals the vanishing order at \(z=0\) of the corresponding theta function with characteristic.

Define
\[
G(g)=\left\{
[A_\tau]\in\mathcal A_g:
\begin{array}{l}
\text{there exist }\varepsilon,\delta\in\{0,1\}^g
\text{ with }\varepsilon^{\mathsf T}\delta\equiv1\pmod2,\\
\operatorname{ord}_{z=0}
\theta\!\begin{bmatrix}\varepsilon\\\delta\end{bmatrix}(\tau,z)\geq3
\end{array}
\right\}.
\]
The use of a symmetric theta representative is part of this interpretation of odd two-torsion; the existential condition is expressed on the moduli space of principally polarized abelian varieties.

\textbf{Conjectural statement supplied by the registry.} For every genus \(g\) in the intended range, is \(G(g)\) pure of codimension \(g\) in \(\mathcal A_g\)? Explicitly, the proposed condition is
\[
\operatorname{codim}_{\mathcal A_g}Z=g
\quad\text{for every irreducible component }Z\text{ of }G(g),
\]
or equivalently
\[
\dim Z=\frac{g(g+1)}2-g=\frac{g(g-1)}2.
\]
The source does not specify how empty loci are to be treated or whether any small genera are excluded. The componentwise formulation alone does not impose nonemptiness. Moreover, the supplied literature assessment says that the underlying excerpt omits the full conjectural claim, despite the registry supplying the displayed formulation and labeling it verified. Thus the intended scope and completeness of this formulation remain unresolved.
\subsection{Short English statement}
The supplied conjecture asks whether, in every intended genus \(g\), every irreducible component of the locus of principally polarized abelian varieties with a theta divisor of multiplicity at least three at some odd two-torsion point has codimension \(g\). The complex setting, genus range, empty-locus convention, and completeness of the extracted claim require confirmation. The supplied open-status label is not independently verified.
\subsection{Sources}
\begin{itemize}
\item Local registry: supplied UnsolvedMath record, permanent ID \(30001576\), title \emph{Triple-Point Theta Divisors at Odd Two-Torsion Points}, source problem code \texttt{catalogue entry}, in the Oberwolfach Reports Open Problems collection. No individual proposer is supplied; the proposed year is recorded as 2010.
\item Supplied report citation: \emph{Komplexe Analysis (2011)}, Oberwolfach Reports, pp.\ 2283--2333, DOI \texttt{10.4171/owr/2010/39}. Original OWR report: \url{https://doi.org/10.4171/owr/2010/39}. The citation year and DOI are retained exactly as supplied.
\item Supplied literature source: Debarre, complex algebraic geometry survey, \url{https://library2.msri.org/books/Book28/files/debarre.pdf}. No fuller bibliographic identification is supplied.
\item Literature assessment dated \texttt{2026-08-21}: supplied dated curation, not independently verified. It labels the problem open, states that the excerpt records low-genus cases, including genus five, and says that these do not establish the general statement. It also states that the excerpt omits the actual full claim and requests independent verification of time-sensitive status and cited claims.
\item The registry labels the statement \texttt{corrected\_verified} and reports that its mathematical content was checked against the official Oberwolfach report, with notation, context, and question form normalized without changing the source-level claim. This is supplied verification metadata, not verification performed for this draft.
\end{itemize}
\end{document}
